% Notes of the Montreal seminar, 1970 (cote 129, typed part): archive cote 129, site page 42 (PDF page 43).
% typescript. Draft compared: none.
% AI-transcribed by Claude Opus 5.5 (claude-opus-5-5, adaptive thinking, effort high), Message Batches API, batch mtl129_full (msgbatch_01AGnDKHwy6ATCM7ehUa6263), 2026-11-13; not yet checked by eye.
\scanpage{129/42}
\mspage{II.14}

\margin{$\times$} \hand{conduisant aux identifications}
\[ D^*(\over{F}{\mathbb{F}}_{G/k}) = F_{D^*(G)} \]
\[ D^*(\over{V}{\mathbb{V}}_G) = V_{D^*(G)} \ . \]

\uline{Corollaire 4.2.1}. On a les équivalences suivantes :

$G$ est infinitésimal $\Longleftrightarrow$ $F$ opère de façon nilpotente sur $D^*(G)$ .

$G$ est unipotent $\Longleftrightarrow$ $V$ opère de façon nilpotente sur $D^*(G)$ .

$G$ est bi-infinitésimal $\Longleftrightarrow$ $F$ et $V$ opèrent de façon nilpotente.

$G$ est étale $\Longleftrightarrow$ $F$ opère bijectivement sur $D^*(G)$ .

$G$ est de type multiplicatif $\Longleftrightarrow$ $V$ opère bijectivement sur $D^*(G)$ .

Pour une démonstration complète du théorème, se reporter
au livre de P. Gabriel et M. Demazure, ``Groupes algébriques''. Dans ce
qui suit, on donnera seulement la construction du foncteur $D^*$ .

On observe tout d'abord que, d'après 3.2., la catégorie
p\_Gr-fin(k) est équivalente à la catégorie produit de la sous-catégo-%
rie pleine de p-Gr-fin(k) dont les objets sont les groupes unipotents
et de la sous-catégorie pleine de p-Gr-fin(k) dont les objets sont
les groupes de type multiplicatif. Il suffira donc de construire $D^*$
séparément sur ces deux sous-catégories.

4.3. \uline{Cas des groupes unipotents}

4.3.1. Pour tout groupe unipotent sur $k$ , on pose par définition
\[ D^*(G) = \mathrm{Hom}_{k\text{-}Gr}(G,\underset{\rightarrow}{\mathbb{W}}) \]
où, pour abréger, on écrit $\underset{\rightarrow}{\mathbb{W}}$ pour $\underset{\rightarrow}{\mathbb{W}}_k$ (I,3). On sait que $W$ agit sur
$\underset{\rightarrow}{\mathbb{W}}$ (I,4.2) et donc $D^*(G)$ est un $W$-module, qui est de longueur finie si
$G$ est fini sur $k$ .

\mspage{42}
